The \code{BPatch\_binaryEdit} class represents a set of executable
files and library files for binary rewriting.
\code{BPatch\_binaryEdit} inherits from the \code{BPatch\_addressSpace}
class, where most functionality for binary rewriting is
found.

\begin{apient}
  bool writeFile(const char *outFile)
\end{apient}

\apidesc{
  Rewrite a \code{BPatch\_binaryEdit} to disk. The original file
  opened with this \code{BPatch\_binaryEdit} is written to the current
  working directory with the name outFile. If any dependent libraries
  were also opened and have instrumentation or
  other modifications, then those libraries will be written to disk
  in the current working directory under their original
  names.

  A rewritten dependency library should only be used with the
  original file that was opened for rewriting. For example,
  if the file a.out and its dependent library libfoo.so were opened
  for rewriting, and both had instrumentation inserted,
  then the rewritten libfoo.so should not be used without the
  rewritten a.out. To build a rewritten libfoo.so that can
  load into any process, libfoo.so must be the original file opened
  by \code{BPatch::openBinary}.

  This function returns true if it successfully wrote a file, or
  false otherwise.
}
